\scanpage{27}
% archive sheet number: 26
% source lines 1-7
Pour modifier cette situation et ouvrir la voie à une politique vraiment nouvelle il faut affaiblir de manière décisive l'emprise politique et économique des monopoles. Toute autre perspective serait illusoire. En refusant les nationalisations indispensables, en proposant de laisser le capital privé maître de ses filiales, en réduisant les mesures fiscales contre les milliardaires, le PS et le MRG vident le programme commun de son efficacité.

% source lines 8-14
Nous partageons la déception et l'inquiétude de tous ceux qui ont placé leur espoir dans le programme commun. Cependant une voie existe pour assurer la victoire de la gauche, non pas seulement la victoire électorale mais le succès de la politique anti-monopoliste à laquelle elle ne peut renoncer sans se désavouer. La gauche française ne doit pas se rabattre sur une politique de gestion de la crise, à la manière des partis socio-démocrates au pouvoir.

% source lines 15-17
La seule voie est la lutte immédiate pour obtenir la signature d'un programme commun actualisé, c'est-à-dire ajusté à une situation très aggravée depuis 1972.

% source lines 18-21
Nous appelons tous les membres du personnel des Universités à prendre leurs responsabilités. Dans ce moment crucial l'avenir dépend de vous. Plus que jamais c'est vous, c'est nous qui pouvons faire changer le Parti Socialiste.

% source lines 22-24
C'est pourquoi nous vous appelons à soutenir nos propositions non parce que ce sont celles du Parti Communiste mais parce qu'elles sont désormais les seules fidèles au programme de 1972 et à sa portée transformatrice.

% source line 25
\begin{center}
POUR LE CHANGEMENT PAR L'UNION SUR UN BON PROGRAMME
\end{center}

% source lines 26-27
\begin{center}
MANIFESTEZ LE 10 DECEMBRE 1977 A 15 H à MONTPELLIER
\end{center}
\begin{flushright}
(rassemblement aux Arceaux)
\end{flushright}

% source line 28
\begin{flushright}
Section des Universités du Parti Communiste Français
\end{flushright}
